% TOP_PROBLEM: {"id": 3228, "title": "Oriented chromatic number of planar graphs", "category_id": 3, "category": {"id": 3, "name": "graph_theory", "display_name": "Graph Theory", "description": "Problems involving graphs, networks, and their properties.", "slug": "graph-theory", "order_index": 3, "created_at": "2026-07-31T15:26:25.671Z"}, "status": "open", "problem_number": "OPG-494", "set_id": 12}
% FIRST_POSTED: 2026-10-01
% ENTRY_KIND: statement_only
% UnsolvedMath ID 3228; source catalogue code OPG-494
% !TeX program = lualatex
% Definition and sources only; this file is not an LLM solution attempt.
\documentclass[11pt,a4paper]{article}
\usepackage[margin=25mm]{geometry}
\usepackage{fontspec}
\setmainfont{lmroman10-regular.otf}[BoldFont=lmroman10-bold.otf,
 ItalicFont=lmroman10-italic.otf,BoldItalicFont=lmroman10-bolditalic.otf]
\usepackage{amsmath,amssymb,mathtools}
\usepackage{unicode-math}
\setmathfont{latinmodern-math.otf}
\AtBeginDocument{\let\setminus\smallsetminus}
\usepackage{xurl,hyperref}
\hypersetup{colorlinks=true,urlcolor=blue}
\setcounter{secnumdepth}{0}
\setlength{\parindent}{0pt}
\setlength{\parskip}{5pt}
\setlength{\emergencystretch}{3em}
\begin{document}
\section*{Oriented chromatic number of planar graphs}\label{3228}
{\small UnsolvedMath ID 3228 · OPG-494 · Graph Theory · OpenGarden}
\subsection{Definitions and mathematical statement}
An oriented graph is obtained from a finite simple undirected graph by choosing exactly one direction for each edge. In particular it has neither loops nor a pair of opposite arcs. An oriented $k$-coloring is a map $c:V\to\{1,\ldots,k\}$ with two properties: endpoints of every arc have different colors, and whenever $u\to v$ and $x\to y$ are arcs with $c(u)=c(y)$, one has $c(v)\ne c(x)$. Thus arcs between any two color classes, if present, all point in the same direction.

The oriented chromatic number $\chi_o(D)$ is the least $k$ admitting such a coloring. Equivalently, it is the least order of an oriented graph receiving a vertex homomorphism from $D$, where a homomorphism sends each arc to an arc of the same direction. The target may be completed to a tournament without harming a homomorphism.

The problem is to determine the exact universal value
\[
 P=\sup\{\chi_o(D):\text{the underlying undirected graph of }D
                         \text{ is planar}\}.
\]
The supremum ranges over every finite planar graph and every orientation of its edges. It is not the minimum over orientations of each graph. A complete answer must supply a uniform upper bound and a matching lower construction, or establish unboundedness if no finite upper bound exists. The target graph for the homomorphism may vary with $D$; a single common tournament is not required by this definition.
\subsection{Short English statement}
Determine the largest oriented chromatic number obtainable by orienting the edges of a finite planar graph.
\subsection{Sources}
\begin{itemize}
\item[S1] ULAM.AI and the UnsolvedMath Contributors, supplied problems.json, record 3228 (OPG-494), OpenGarden. Catalogue identity and source transcription; CC BY 4.0.
\url{https://huggingface.co/datasets/ulamai/UnsolvedMath}.
\item[S2] Open Problem Garden, Oriented chromatic number of planar graphs. Original problem and discussion.
\url{http://www.openproblemgarden.org/op/oriented_chromatic_number_of_planar_graphs}.
\end{itemize}
\end{document}
